% Generated draft: requires mathematical and source review.
\documentclass[11pt]{article}
\usepackage[margin=25mm]{geometry}
\usepackage{fontspec}
\setmainfont{Latin Modern Roman}
\usepackage{amsmath,amssymb,mathtools,unicode-math}
\setmathfont{Latin Modern Math}
\usepackage{enumitem,needspace,xurl,hyperref}
\setlength{\parindent}{0pt}
\setlength{\parskip}{5pt}
\setlength{\emergencystretch}{3em}
\newcommand{\N}{\mathbb{N}}
\newcommand{\Z}{\mathbb{Z}}
\newcommand{\Q}{\mathbb{Q}}
\newcommand{\R}{\mathbb{R}}
\newcommand{\C}{\mathbb{C}}
\newcommand{\F}{\mathbb{F}}
\newcommand{\A}{\mathbb{A}}
\newcommand{\PP}{\mathbb{P}}
\newcommand{\E}{\mathbb{E}}
\newcommand{\eps}{\varepsilon}
\newcommand{\dd}{\,\mathrm{d}}
\begin{document}
\section*{UnsolvedMath ID: 30000323}
\textbf{Source problem code:} catalogue entry
\subsection{Definitions and mathematical statement}
Let $p$ be an odd prime, $G$ a finite abelian $p$-group, and $V$ a finite-dimensional complex $G$-representation with no trivial direct summand. Equivalently, $V^G=0$. Since $G$ is abelian, it decomposes as a sum of nontrivial one-dimensional characters; faithfulness of the representation is not assumed, and $V^G=0$ does not assert that every nonidentity element acts without fixed vectors.

Let $BG=EG/G$ and let $\xi_V=EG\times_G V\to BG$ be the associated bundle, of complex rank $r$. In an oriented bordism cohomology theory $E$ in which this bundle has its specified Thom class $u_E(\xi_V)$, its Euler class is
\[
 e_E(V)=s^*u_E(\xi_V)\in E^{2r}(BG),
\]
where $s$ is the zero section. The question is whether
\[
 e_E(V)\ne0
\]
in the topological bordism theory intended by the report, for every such $G$ and $V$.

The extract does not name the representing bordism spectrum, orientation convention or localization/completion. Oriented topological bordism, complex cobordism and unoriented bordism are different theories; one cannot replace the intended theory by ordinary cohomology or by a chosen alternative simply to interpret the Euler class. The displayed $E$ notation exposes this missing specification. If an orientation supporting the Whitney-sum formula is fixed, the character decomposition expresses the Euler class as the product of the corresponding line-bundle Euler classes, but nonvanishing of the factors alone does not prove nonvanishing of their product.

The literature assessment checked on 21 August 2026 says the fixed-point transfer theorem in the source passage does not answer this Euler-class question and treats it as unresolved. No stronger consequence of that transfer theorem is supplied. This is dated triage requiring independent verification.
\subsection{Short English statement}
For an odd-primary finite abelian group, must a complex representation with no trivial summand have nonzero Euler class in the report's topological bordism theory? The exact bordism theory and orientation must be specified; fixed-point transfer alone is not the supplied answer.
\subsection{Sources}
\begin{itemize}
\item \emph{Cohomology of Finite Groups: Interactions and Applications} (2006), Oberwolfach Reports, DOI 10.4171/owr/2005/42, pp. 2375--2434, \url{https://doi.org/10.4171/owr/2005/42}; proposed year supplied as 2005.
\item The sole literature source in the assessment checked 21 August 2026 is the original report, \url{https://doi.org/10.4171/owr/2005/42}. It distinguishes the Euler-class question from the fixed-point transfer theorem in the passage.
\item Local UnsolvedMath registry, permanent ID $30000323$, exact source code \texttt{catalogue entry}, created 25 August 2026 and updated 16 September 2026. The clean statement is labelled corrected and verified, restoring the finite abelian odd-$p$ group and representation hypotheses.
\end{itemize}
\end{document}
